% =============================================================================
%  Chapter 59 — Incremental Control Policy for NonlinearAnalysis
% =============================================================================

\chapter{Incremental Control Policy for Nonlinear Analysis}
\label{ch:incremental-control}

This chapter documents the design, implementation, and testing of the
\texttt{IncrementalControlPolicy} abstraction, which generalises the
incremental stepping strategy of \texttt{NonlinearAnalysis} from
hardcoded proportional load control to an injectable, compile-time
policy.  Three concrete schemes are provided: \texttt{LoadControl}
(backward-compatible default), \texttt{DisplacementControl} (single
or multi-DOF pushover), and \texttt{CustomControl<F>} (user-supplied
callable for maximum flexibility).  An \texttt{ArcLengthControl} stub
is included for future development.

% =============================================================================
\section{Motivation and Problem Statement}
\label{sec:ic-motivation}

The original implementation of \texttt{NonlinearAnalysis::advance\_to\_lambda\_}
contained three hardcoded lines that coupled the class to proportional load
control:

\begin{lstlisting}
VecCopy(f_full, f_ext);
VecScale(f_ext, lambda_target);
ctx_.f_ext = f_ext;
\end{lstlisting}

This prevented using the existing bisection infrastructure for:

\begin{itemize}
  \item \textbf{Displacement control} — prescribing displacement increments
        at specific DOFs (pushover analysis, post-peak softening).
  \item \textbf{Mixed control} — combining load scaling with prescribed
        displacements.
  \item \textbf{Custom ramps} — non-linear load paths, conditional logic,
        or external coupling.
  \item \textbf{Arc-length control} (future) — parametrising the solution
        path by its arc-length to traverse limit points.
\end{itemize}

The goal was to extract this logic into a \emph{compile-time injectable
policy} that:
\begin{enumerate}[label=(\roman*)]
  \item preserves zero overhead (monomorphised at compile time),
  \item is backward-compatible (existing code compiles without changes),
  \item enables displacement control, custom ramps, and future arc-length,
  \item integrates seamlessly with the existing bisection engine.
\end{enumerate}

% =============================================================================
\section{Design Decision: Function-Level Template}
\label{sec:ic-design}

Two injection points were considered:

\begin{description}
  \item[Class-level template parameter] Add \texttt{CS} to the class template
        parameter list of \texttt{NonlinearAnalysis<\ldots, CS>}.  This forces
        the control scheme to be fixed at object construction time and pollutes
        the type name.

  \item[Function-level template parameter] Make \texttt{solve\_incremental}
        a member function template accepting the scheme by value:
        \begin{lstlisting}
template <typename CS>
bool solve_incremental(int num_steps, int max_bisections, CS scheme);
        \end{lstlisting}
        The type \texttt{CS} is deduced from the argument via CTAD; a
        default overload taking no scheme calls with \texttt{LoadControl\{\}}.
\end{description}

The function-level approach was chosen because:

\begin{itemize}
  \item The same \texttt{NonlinearAnalysis} object can solve with
        different schemes in sequence (e.g., load-controlled pre-load
        followed by displacement-controlled pushover).
  \item No additional member state is introduced: the scheme is a
        stack-local value, passed through the recursion.
  \item The default overload \texttt{solve\_incremental(num\_steps)} is
        100\% backward-compatible — no existing call site needs changes.
  \item CTAD deduces \texttt{CS}, so client code never spells the type:
        \begin{lstlisting}
nl.solve_incremental(20, 4, DisplacementControl{42, 0, 0.01});
        \end{lstlisting}
\end{itemize}

This follows the same philosophy as the project's existing policy-based
patterns (\texttt{MaterialPolicy}, \texttt{KinematicPolicy},
\texttt{LocalStepControlPolicy}), but at the function level rather than
the class level, offering greater flexibility without additional
coupling.

% =============================================================================
\section{The \texttt{IncrementalControlPolicy} Concept}
\label{sec:ic-concept}

The concept constrains any type \texttt{S} that can set the system
state for a scalar control parameter \texttt{p}:

\begin{lstlisting}
template <typename S, typename ModelT>
concept IncrementalControlPolicy = requires(
    S& scheme, double p, Vec f_full, Vec f_ext, ModelT* model)
{
    scheme.apply(p, f_full, f_ext, model);
};
\end{lstlisting}

The single required expression \texttt{scheme.apply(p, f\_full, f\_ext,
model)} sets the complete system state for the control parameter
\texttt{p}:

\begin{itemize}
  \item \texttt{p} --- control parameter in $[0, 1]$.  At $p=0$ the
        system is unloaded; at $p=1$ the full reference load or target
        displacement is reached.
  \item \texttt{f\_full} --- reference external force vector (read-only,
        the unscaled load at $p=1$).
  \item \texttt{f\_ext} --- external force vector to fill.  Consumed by
        the SNES residual callback.
  \item \texttt{model} --- pointer to the \texttt{Model} object, giving
        access to \texttt{update\_imposed\_value()}, the constraint map,
        etc.
\end{itemize}

\subsection{The Absolute-State Invariant}
\label{ssec:ic-absolute}

\begin{remark}[Absolute invariant]
\label{rem:ic-absolute}
\texttt{apply(p, \ldots)} must be \emph{absolute}: for any value of
\texttt{p}, calling \texttt{apply(p, \ldots)} fully determines the
system state, regardless of the previous call history.
This is essential for correctness of the bisection engine.
\end{remark}

When a Newton step diverges, the bisection engine:
\begin{enumerate}
  \item Reverts the displacement field $\mathbf{u}$ to its snapshot.
  \item Reverts the material state to the last committed values.
  \item Calls \texttt{scheme.apply(p\_current, \ldots)} to restore
        the control state.
\end{enumerate}

Because \texttt{apply()} is absolute, step~3 requires no save/restore
logic inside the scheme itself --- the system state is a pure function
$f(p)$ of the control parameter.  This eliminates an entire class of
state management bugs and keeps the scheme implementations stateless
and trivially correct.

% =============================================================================
\section{Concrete Implementations}
\label{sec:ic-implementations}

\subsection{\texttt{LoadControl} --- Proportional Load Scaling}
\label{ssec:ic-load}

The default scheme extracts the original three-line logic:

\begin{lstlisting}
struct LoadControl {
    template <typename ModelT>
    void apply(double p, Vec f_full, Vec f_ext,
               [[maybe_unused]] ModelT*) const
    {
        VecCopy(f_full, f_ext);
        VecScale(f_ext, p);
    }
};
\end{lstlisting}

At each step $k$ of $N$, the external force vector is set to
$\mathbf{f}_\text{ext} = p \cdot \mathbf{f}_\text{full}$ where
$p = k/N$.  The model pointer is unused.  This is the identity
refactoring of the original code and serves as the default:

\begin{lstlisting}
nl.solve_incremental(10);           // LoadControl{} implicit
nl.solve_incremental(10, 4, LoadControl{});  // explicit
\end{lstlisting}

\subsection{\texttt{DisplacementControl} --- Prescribed-DOF Pushover}
\label{ssec:ic-displacement}

Controls one or more DOFs linearly with $p$, with optional proportional
load scaling:

\begin{lstlisting}
struct DisplacementControl {
    struct PrescribedDOF {
        std::size_t node;
        std::size_t dof;
        double      target;
    };

    std::vector<PrescribedDOF> controlled_dofs;
    double load_scale{0.0};

    // Single DOF
    DisplacementControl(std::size_t node, std::size_t dof,
                        double target, double scale = 0.0);

    // Multi-DOF
    DisplacementControl(std::vector<PrescribedDOF> dofs,
                        double scale = 0.0);

    template <typename ModelT>
    void apply(double p, Vec f_full, Vec f_ext, ModelT* model) {
        VecCopy(f_full, f_ext);
        VecScale(f_ext, p * load_scale);
        for (const auto& [node, dof, target] : controlled_dofs) {
            model->update_imposed_value(node, dof, p * target);
        }
    }
};
\end{lstlisting}

At control parameter $p$:
\begin{align}
  u_{\text{imposed}, i} &= p \cdot \mathtt{target}_i, \label{eq:ic-disp-dof} \\
  \mathbf{f}_\text{ext} &= p \cdot \mathtt{load\_scale} \cdot \mathbf{f}_\text{full}.
  \label{eq:ic-disp-fext}
\end{align}

The controlled DOFs must have been pre-constrained via
\texttt{constrain\_dof()} before \texttt{Model::setup()}, so that the DM
section includes the boundary condition.  The \texttt{apply()} call only
updates the \emph{value} of an existing constraint (not the section
structure), using \texttt{Model::update\_imposed\_value()}.

Typical usage for a pushover analysis at node~42, DOF~0, with
target displacement 0.01:

\begin{lstlisting}
M.constrain_dof(42, 0, 0.0);   // register constraint before setup
M.setup();

NonlinearAnalysis<ThreeDimensionalMaterial> nl{&M};
nl.solve_incremental(20, 4, DisplacementControl{42, 0, 0.01});
\end{lstlisting}

\subsection{\texttt{CustomControl<F>} --- Maximum Flexibility}
\label{ssec:ic-custom}

Wraps any callable with signature
\texttt{void(double, Vec, Vec, ModelT*)} into a policy-satisfying
object:

\begin{lstlisting}
template <typename F>
struct CustomControl {
    F apply_fn;

    template <typename ModelT>
    void apply(double p, Vec f_full, Vec f_ext, ModelT* model) {
        apply_fn(p, f_full, f_ext, model);
    }
};

template <typename F> CustomControl(F) -> CustomControl<F>;

template <typename F>
auto make_control(F&& f) {
    return CustomControl<std::decay_t<F>>{std::forward<F>(f)};
}
\end{lstlisting}

This enables arbitrary control strategies without defining a named
class.  Example with a square-root ramp combined with displacement
control:

\begin{lstlisting}
nl.solve_incremental(10, 4, make_control(
    [](double p, Vec f_full, Vec f_ext, auto* model) {
        VecCopy(f_full, f_ext);
        VecScale(f_ext, std::sqrt(p));
        model->update_imposed_value(42, 0, p * 0.01);
    }
));
\end{lstlisting}

The CTAD deduction guide and \texttt{make\_control()} factory work
in tandem: the factory perfectly forwards the callable and strips
references, while the deduction guide enables direct aggregate
initialisation when desired.

\subsection{Arc-Length Control (Stub)}
\label{ssec:ic-arclength}

An \texttt{ArcLengthControl} stub is placed in
\texttt{IncrementalControl.hh} for future development.  Arc-length
(Riks) methods require augmenting the SNES system to the unknowns
$(\mathbf{u}, \lambda)$ with a constraint equation:
\[
  \|\Delta\mathbf{u}\|^2 + \Delta\lambda^2 \|\mathbf{f}\|^2 = \Delta s^2,
\]
where $s$ is the arc-length parameter.  The control parameter becomes
$s$ rather than $p$, and the stepping engine must reinterpret its
meaning accordingly.  This is marked with a \texttt{static\_assert}
guard preventing accidental instantiation.

% =============================================================================
\section{\texttt{Model::update\_imposed\_value()}}
\label{sec:ic-update-imposed}

The \texttt{DisplacementControl} scheme needs to update the imposed value
of a constraint at each step without rebuilding the DM section or
retraversing all nodes.  The existing \texttt{fill\_imposed\_solution()}
method is $O(N)$ --- it walks every node in the domain --- which is
unacceptable inside the stepping loop.

A new $O(1)$ method was added to \texttt{Model}:

\begin{lstlisting}
void update_imposed_value(std::size_t node_idx,
                          std::size_t local_dof,
                          double value) noexcept
{
    // 1. Update in-memory constraint map (O(log n) amortised)
    constrain_dof(node_idx, local_dof, value);

    // 2. Direct write to PETSc imposed-solution Vec (O(1))
    auto& node = domain_->node(node_idx);
    auto  all_dofs = node.dof_index();
    PetscScalar val = static_cast<PetscScalar>(value);
    VecSetValueLocal(global_imposed_solution_.get(),
                     all_dofs[local_dof], val, INSERT_VALUES);
    VecAssemblyBegin(global_imposed_solution_.get());
    VecAssemblyEnd(global_imposed_solution_.get());
}
\end{lstlisting}

This method:
\begin{enumerate}
  \item Calls \texttt{constrain\_dof()} to update the constraint map
        (the same map used by \texttt{fill\_imposed\_solution()}).
  \item Writes the new value directly into the PETSc Vec via
        \texttt{VecSetValueLocal()}, using the DOF's local index
        from the node. This avoids the full-domain traversal.
  \item Calls \texttt{VecAssemblyBegin/End} to flush the value
        (required by PETSc even for sequential vectors).
\end{enumerate}

\begin{remark}
The DOF must have been registered as constrained \emph{before}
\texttt{setup()} so that the DM section includes the boundary
condition.  \texttt{update\_imposed\_value()} only updates the value,
not the topological structure.
\end{remark}

% =============================================================================
\section{Changes to \texttt{NonlinearAnalysis}}
\label{sec:ic-nlanalysis}

\subsection{Renamed Control Parameter}
\label{ssec:ic-rename}

The load factor $\lambda$ was renamed to the generic control parameter
$p$ throughout the implementation.  The stepping function
\texttt{advance\_to\_lambda\_} became \texttt{advance\_to\_p\_}.
Print messages were updated accordingly.  This reflects the fact that
the control parameter may represent load scaling, displacement fraction,
or any other monotonic path parameter.

\subsection{Template Structure}
\label{ssec:ic-template}

The two modified functions are now templates on the control scheme:

\begin{lstlisting}
// Backward-compatible overload
bool solve_incremental(int num_steps, int max_bisections = 4);

// General form
template <typename CS>
bool solve_incremental(int num_steps, int max_bisections, CS scheme);

// Bisection engine (private)
template <typename CS>
bool advance_to_p_(CS& scheme, Vec f_full,
                   double p_current, double p_target,
                   int bisections_left);
\end{lstlisting}

The backward-compatible overload delegates to the general form with
\texttt{LoadControl\{\}} as the default scheme.  A \texttt{static\_assert}
in the general form enforces the concept:

\begin{lstlisting}
static_assert(IncrementalControlPolicy<CS, ModelT>,
    "CS must satisfy IncrementalControlPolicy<CS, ModelT>");
\end{lstlisting}

\subsection{Integration with the Bisection Engine}
\label{ssec:ic-bisection}

The bisection engine (\texttt{advance\_to\_p\_}) received minimal
changes.  The only modification is the replacement of the three
hardcoded lines with a single call:

\begin{lstlisting}
// Before (hardcoded):
VecCopy(f_full, f_ext);
VecScale(f_ext, lambda_target);
ctx_.f_ext = f_ext;

// After (policy-driven):
scheme.apply(p_target, f_full, f_ext, model_);
ctx_.f_ext = f_ext;
\end{lstlisting}

The snapshot/revert logic for the displacement field $\mathbf{u}$ and
material state is unchanged.  After a diverged step, the engine calls
\texttt{scheme.apply(p\_current, \ldots)} to restore the control
state before bisecting --- this is correct because \texttt{apply()}
is absolute (see Remark~\ref{rem:ic-absolute}).

The full bisection algorithm is summarized in
Algorithm~\ref{alg:ic-bisection}.

\begin{algorithm}[ht]
  \caption{Policy-driven bisection engine (\texttt{advance\_to\_p\_}).}
  \label{alg:ic-bisection}
  \SetAlgoLined
  \KwIn{scheme, $\mathbf{f}_\text{full}$, $p_\text{current}$,
        $p_\text{target}$, budget}
  \KwOut{\texttt{true} if $p_\text{target}$ reached,
         \texttt{false} otherwise}

  $\mathbf{u}_\text{snap} \gets \mathbf{u}$\;
  \texttt{scheme.apply}$(p_\text{target})$\;
  Solve SNES\;
  \uIf{converged}{
    Commit material state\;
    \Return \texttt{true}\;
  }
  \Else{
    $\mathbf{u} \gets \mathbf{u}_\text{snap}$\;
    Revert material state\;
    \texttt{scheme.apply}$(p_\text{current})$ \tcp*{restore control state}
    \uIf{budget $= 0$}{
      \Return \texttt{false}\;
    }
    \Else{
      $p_\text{mid} \gets (p_\text{current} + p_\text{target}) / 2$\;
      \lIf{\textbf{not} advance\_to\_p\_(scheme, $p_\text{current}$,
            $p_\text{mid}$, budget$-1$)}{\Return \texttt{false}}
      \Return advance\_to\_p\_(scheme, $p_\text{mid}$,
              $p_\text{target}$, budget$-1$)\;
    }
  }
\end{algorithm}

% =============================================================================
\section{Relationship with Existing Policy Patterns}
\label{sec:ic-patterns}

The \texttt{IncrementalControlPolicy} fits naturally into the
project's policy-based architecture.  \autoref{tab:ic-policy-summary}
contextualises it alongside existing abstractions.

\begin{table}[ht]
  \centering
  \caption{Policy abstractions in \texttt{fall\_n} at different scales.}
  \label{tab:ic-policy-summary}
  \begin{tabular}{lll}
    \toprule
    Scale & Concept / Interface & Implementations \\
    \midrule
    Material (local) & \texttt{LocalStepControlPolicy}
      & \texttt{FullStepControl},
        \texttt{BacktrackingResidualStepControl} \\
    Dynamic analysis & \texttt{ForceEvaluator} (\texttt{std::function})
      & Lambdas, closures \\
    \textbf{Static NL}
      & \textbf{\texttt{IncrementalControlPolicy}}
      & \textbf{\texttt{LoadControl}},
        \textbf{\texttt{DisplacementControl}}, \\
      & & \textbf{\texttt{CustomControl<F>}} \\
    \bottomrule
  \end{tabular}
\end{table}

The key difference from \texttt{LocalStepControlPolicy} (which also
uses C++20 concepts) is the injection point: the local step control
is a class-level template parameter of the material, while the
incremental control is a function-level template on
\texttt{solve\_incremental}.  This reflects their different scopes:
material policies are fixed for the lifetime of an element, while
the incremental control strategy may change between analysis phases.

The \texttt{ForceEvaluator} pattern used in
\texttt{DynamicAnalysis} relies on \texttt{std::function} for
type-erasure (suitable for time-dependent forces that vary at
runtime).  In contrast, \texttt{IncrementalControlPolicy} uses
compile-time polymorphism because the control scheme is typically
fixed for an entire \texttt{solve\_incremental} call and
performance in the inner loop is critical.

% =============================================================================
\section{File Summary}
\label{sec:ic-files}

\autoref{tab:ic-files} summarises all changes made.

\begin{table}[ht]
  \centering
  \caption{Files created or modified for the incremental control policy.}
  \label{tab:ic-files}
  \begin{tabular}{lll}
    \toprule
    File & Type & Description \\
    \midrule
    \texttt{src/analysis/IncrementalControl.hh}
      & New & Concept + \texttt{LoadControl} +
              \texttt{DisplacementControl} +
              \texttt{CustomControl<F>} \\
    \texttt{src/analysis/NLAnalysis.hh}
      & Modified & Template \texttt{solve\_incremental} +
                   \texttt{advance\_to\_p\_} \\
    \texttt{src/model/Model.hh}
      & Modified & Added \texttt{update\_imposed\_value()} \\
    \texttt{src/analysis/Analysis.hh}
      & Modified & Added include for \texttt{IncrementalControl.hh} \\
    \texttt{tests/test\_incremental\_control.cpp}
      & New & 6~tests, 14~assertions \\
    \texttt{CMakeLists.txt}
      & Modified & Registered test target \\
    \bottomrule
  \end{tabular}
\end{table}

% =============================================================================
\section{Testing}
\label{sec:ic-testing}

Six unit tests validate the implementation across all control schemes.
The test file \texttt{tests/test\_incremental\_control.cpp} creates a
minimal 3D unit-cube model (single hexahedral element, elastic material)
and exercises:

\begin{enumerate}
  \item \textbf{LoadControl (default call)} ---
        \texttt{solve\_incremental(3)} with no explicit scheme;
        verifies backward compatibility and convergence in 3~steps.

  \item \textbf{LoadControl (explicit)} ---
        \texttt{solve\_incremental(3, 4, LoadControl\{\})} produces
        the same result.

  \item \textbf{DisplacementControl (single DOF)} ---
        Constrains node~1 $x$-DOF with target~$0.01$; verifies
        convergence and that the imposed solution vector is non-zero
        at the controlled DOF.

  \item \textbf{DisplacementControl (multi-DOF)} ---
        Four nodes controlled simultaneously; verifies convergence.

  \item \textbf{CustomControl via \texttt{make\_control}} ---
        Lambda with $\sqrt{p}$ load ramp; verifies convergence.

  \item \textbf{Concept satisfaction (compile-time)} ---
        Static assertions that \texttt{LoadControl},
        \texttt{DisplacementControl}, and
        \texttt{CustomControl<\ldots>} satisfy
        \texttt{IncrementalControlPolicy}.
\end{enumerate}

All 14~assertions pass.  The existing 30~commit-chain tests and the
heterogeneity test (which uses \texttt{solve\_incremental}) also pass
without modification, confirming backward compatibility.

% =============================================================================
\section{Usage Guide}
\label{sec:ic-usage}

\subsection{Load Control (Default)}

No changes to existing code:

\begin{lstlisting}
NonlinearAnalysis<ThreeDimensionalMaterial> nl{&M};
nl.solve_incremental(10);        // LoadControl implicit, 10 steps
nl.solve_incremental(10, 4);     // same with explicit bisection depth
\end{lstlisting}

\subsection{Displacement-Controlled Pushover}

\begin{lstlisting}
// Pre-constrain the controlled DOF (zero initial value)
M.constrain_dof(42, 0, 0.0);
M.setup();

NonlinearAnalysis<ThreeDimensionalMaterial> nl{&M};
// Push node 42, dof 0 to displacement 0.01 in 20 steps
nl.solve_incremental(20, 4, DisplacementControl{42, 0, 0.01});
\end{lstlisting}

\subsection{Multi-DOF Displacement Control with Proportional Load}

\begin{lstlisting}
nl.solve_incremental(20, 4, DisplacementControl{
    {{10, 0, 0.01}, {11, 0, 0.01}, {12, 0, 0.01}},
    0.5   // load_scale: 50% of reference load
});
\end{lstlisting}

\subsection{Custom Control}

\begin{lstlisting}
nl.solve_incremental(10, 4, make_control(
    [](double p, Vec f_full, Vec f_ext, auto* model) {
        VecCopy(f_full, f_ext);
        VecScale(f_ext, std::sqrt(p));   // sqrt ramp
        model->update_imposed_value(42, 0, p * 0.01);
    }
));
\end{lstlisting}

\subsection{Sequential Multi-Phase Analysis}

Because the scheme is a function-level parameter, the same
\texttt{NonlinearAnalysis} object can be reused with different
schemes:

\begin{lstlisting}
NonlinearAnalysis<ThreeDimensionalMaterial> nl{&M};

// Phase 1: gravity load (proportional)
nl.solve_incremental(5, 4, LoadControl{});

// Phase 2: displacement-controlled pushover
nl.solve_incremental(20, 4, DisplacementControl{42, 0, 0.05});
\end{lstlisting}

This is a direct benefit of the function-level template design
over a class-level template parameter.
